% Templat TA TIF UIN Suska Riau -- (c) 2026 Benny Sukma Negara
\bagianawal{ABSTRAK}

Abstrak merupakan deskripsi singkat keseluruhan proses \tampil{jenis} yang dilakukan. Pada umumnya berisi tentang motivasi penelitian, rumusan masalah, kajian dan eksperimen yang telah dilakukan serta hasil kajian dan eksperimen. Abstrak dibuat singkat, minimal 100 kata dan tidak lebih dari 200 kata, dalam satu paragraf, tanpa gambar, tabel, rumus, nomor urut, maupun subjudul. Kata kunci berjumlah 3--5, diurutkan menurut abjad. \isi{Ganti seluruh paragraf ini dengan abstrak Anda.}

Kata kunci: abstrak, buku panduan, tugas akhir

\bagianawal{ABSTRACT}

The abstract is a brief description of the entire process of the Final Project. In general, it contains research motivation, problem formulation, studies and experiments that have been carried out as well as the results of studies and experiments. \isi{Replace this paragraph with your English abstract.}

Keywords: abstract, final project, manual book
